\documentclass[10pt]{article}
\usepackage[utf8]{inputenc}
\usepackage[T1]{fontenc}
\usepackage[spanish]{babel}
\usepackage[margin=1.8cm]{geometry}
\usepackage{titlesec}
\titlespacing*{\section}{0pt}{1.1ex}{0.5ex}
\titlespacing*{\subsection}{0pt}{0.8ex}{0.3ex}
\usepackage{amsmath,amssymb}
\usepackage{lmodern}
\usepackage{parskip}
\begin{document}
\begin{center}
{\Large\bfseries Probabilidad en Python - Resumen (Parte 2)}
\end{center}
\section*{1. Conceptos fundamentales}
\begin{itemize}
\item Evento: subconjunto del espacio muestral $\Omega$.
\item Asignación de probabilidad: a cada evento A contenido en $\Omega$ se le asigna un número entre 0 y 1 mediante la regla clásica: P(A) = |A|/|$\Omega$|.
\end{itemize}
\section*{2. Definición de la función de probabilidad en Python}
Para calcular automáticamente la probabilidad de un evento, se puede definir una función usando len(), que obtiene la cardinalidad de un conjunto:

\begin{verbatim}
def P(A, Omega):
    return len(A) / len(Omega)
\end{verbatim}
\begin{itemize}
\item len(A): cardinalidad del evento A, es decir, número de casos favorables.
\item len(Omega): cardinalidad del espacio muestral $\Omega$, es decir, número de casos totales.
\end{itemize}
\section*{3. Eventos y operaciones de conjuntos en Python}
\begin{itemize}
\item Unión de eventos (A $\cup$ B): sucede A o sucede B.
\item Intersección de eventos (A $\cap$ B): suceden A y B al mismo tiempo.
\item Diferencia de eventos (A \textbackslash{} B): sucede A pero no sucede B.
\item Complemento de un evento (A$^c$): elementos de $\Omega$ que no pertenecen a A.
\end{itemize}
\section*{4. Probabilidad condicionada}
La probabilidad de que suceda A dado que ya sucedió B se define como:

P(A | B) = P(A $\cap$ B) / P(B), siempre que P(B) > 0.

El archivo original termina indicando que la implementación en Python se realiza con una función que recibe los eventos A y B dentro del espacio muestral $\Omega$; no incluye el código de esa función.

\end{document}